
\subsection{原子模型}
\begin{tabularx}{\columnwidth}{XX}
  卢瑟福散射角判定&大角散射对应近核强库仑作用\\
  玻尔能级公式（氢原子）&$E_n=-\frac{13.6\ \mathrm{eV}}{n^2}$\\
  轨道半径（氢原子）&$r_n=n^2a_0$，$a_0\approx 0.53\times 10^{-10}\ \mathrm{m}$\\
  类氢离子轨道半径&$r_n=\frac{n^2}{Z}a_0$\\
  频率条件&$h\nu=|E_m-E_n|$\\
  最大谱线数（从第 $n$ 能级向下）&$N=\frac{n(n-1)}{2}$\\
\end{tabularx}
\textbf{枣糕模型}
\paragraph{阴极射线} 发现原子可再分，揭示电子存在。

汤姆孙在研究阴极射线时确认了电子的存在，由此说明原子并非不可再分的实心小球。在这个基础上提出“枣糕模型”：原子是带正电的均匀球体，电子镶嵌其中以保证整体电中性。
\textbf{核式结构模型}

卢瑟福通过 $\alpha$ 粒子散射实验发现：绝大多数粒子几乎不偏转直接穿过金箔，只有极少数粒子发生大角度偏转甚至反弹。由此提出核式结构模型：原子内部绝大部分空间是空的，中心存在体积极小但质量很大的原子核，原子核集中几乎全部正电荷和绝大部分质量。

\textbf{波尔量子化轨道模型}

波尔在核式结构基础上引入“轨道量子化”假设：电子只能稳定地处在一组离散定态轨道（能级）上，在定态轨道运动时不辐射能量。

\textbf{能级}：氢原子第 $n$ 能级为
\[
  E_n=-\frac{13.6\ \mathrm{eV}}{n^2}.
\]
$n$ 越大，能级越高（数值上越接近 0），电子束缚越弱。

\textbf{轨道半径}：氢原子第 $n$ 条定态轨道半径为
\[
  r_n=n^2a_0.
\]
其中 $a_0$ 为第一轨道半径。类氢离子（核电荷数为 $Z$，核外只有一个电子）常用
\[
  r_n=\frac{n^2}{Z}a_0.
\]
同一原子中 $n$ 越大，轨道半径越大，相邻轨道半径并不是等差增加。

\textbf{跃迁规律}：
\begin{itemize}
  \item 电子从高能级跃迁到低能级（$n_i>n_f$）时，释放光子，满足
        \[
          h\nu=E_{n_i}-E_{n_f}.
        \]
  \item 电子从低能级跃迁到高能级（$n_i<n_f$）时，必须吸收能量恰好等于能级差，满足
        \[
          h\nu=E_{n_f}-E_{n_i}.
        \]
  \item 当电子获得能量不小于电离能时可脱离原子，终态对应 $E=0$（电离态）。
\end{itemize}

\textbf{跃迁情形判定要点}：先判断“升能级还是降能级”，再判断“吸收还是辐射”；频率和波长由能级差唯一决定，不由轨道编号本身单独决定。

\textbf{谱线所属区域}：氢原子向 $n=1$ 能级跃迁形成莱曼系，光子能量较大，属于紫外线；向 $n=2$ 能级跃迁形成巴耳末系，其中部分谱线位于可见光区域；向 $n=3$ 能级跃迁形成帕邢系，属于红外线。

\begin{formulanotebox}
玻尔模型的能级、半径公式主要适用于氢原子及类氢体系，不能直接机械套用到任意多电子原子。

跃迁辐射频率由能级差决定，与电子“所在轨道编号大小”无直接线性关系。

若题目说“大量氢原子处于第 $n$ 能级”，可出现所有低能级之间的辐射跃迁；若只讨论“一个电子一次从第 $n$ 能级回到基态”，一次只发出一条光谱线，但可能路径不同。
\end{formulanotebox}

\begin{keypointbox}
原子模型演化体现“实验推动理论修正”：从阴极射线发现电子（汤姆孙）到 $\alpha$ 散射建立核式结构（卢瑟福），再到离散能级与跃迁规律（波尔）。
\end{keypointbox}

\begin{casetypebox}[title={\strut\textcolor{white}{\bfseries 常见题型：波尔模型与能级跃迁}}]
\begin{itemize}[leftmargin=1.5em]
  \item \textbf{能级计算}：先写出初、末能级能量，再用 $\Delta E=|E_m-E_n|$，由 $h\nu=\Delta E$ 或 $c=\lambda\nu$ 求频率、波长。
  \item \textbf{吸收与发射判定}：低能级到高能级为吸收，高能级到低能级为发射；吸收光子时能量必须恰好等于两个定态能级差。
  \item \textbf{谱线条数}：大量氢原子从第 $n$ 能级向下跃迁，最多产生 $\frac{n(n-1)}{2}$ 条谱线。
  \item \textbf{第 4 能级向下跃迁}：共有 $4\to3$、$4\to2$、$4\to1$、$3\to2$、$3\to1$、$2\to1$ 六条谱线，其中 $4\to1$、$3\to1$、$2\to1$ 为紫外线，$4\to2$、$3\to2$ 为可见光，$4\to3$ 为红外线。
  \item \textbf{轨道半径判断}：氢原子用 $r_n=n^2a_0$，类氢离子用 $r_n=\frac{n^2}{Z}a_0$；半径比可直接按 $n^2$ 或 $\frac{n^2}{Z}$ 比较。
\end{itemize}
\end{casetypebox}